% STATUS: verified
\chapter{Deployment and operations}\label{ch:ops}

This chapter is for whoever installs the \product{} and keeps it running: the operator. It
follows the order in which an operator meets the system --- what runs and what the host must
offer, a fresh installation, the two modes of the compose file, updates --- and then
turns to the work that is particular to this product: keeping the legal corpus current,
starting the three paid readings of the whole corpus, which only an operator starts, looking after the
prompts, and watching the budgets. It ends with the scheduled work, the health probe and the
common failures. The commands are quoted exactly and kept few; every one runs in the project
folder on the host, and nothing but Docker is installed there. The settings and constants the
chapter names are listed in \cref{app:config}.

Two repository documents cover the same ground in more detail and were the map for this
chapter: \file{docs/markdown/en/deployment.md} and \file{docs/markdown/en/operations.md}. The scripts and the
compose file are the last word, and the chapter was checked against them.


\section{What runs and what the host needs}\label{sec:ops:runs}

The whole system is one Docker Compose project described by \file{docker-compose.yml}: eight
containers on one host (\cref{tab:ops:services,fig:ops:host}). \Cref{ch:arch} explains what the
services do for each other; here they are listed as an operator sees them in
\texttt{docker compose ps}.

Every container has a read-only root filesystem and the restart policy \code{unless-stopped}.
The endpoint, the exit, the API and the database run as user 1000; the console and the
assistant service run as their image's default user (the line \texttt{user: 1000:1000} is
commented out for them); the database tool and the worker container run as root. The worker
container needs root because it holds the host's Docker socket --- which makes it root on the
host --- and for that reason it has no network at all and runs nothing but that one command.
The compose file sets no resource limits, no health checks and no start order: the services
start together, and the endpoint answers 502 until the service behind a path is up.

\begin{table}[H]
\caption{The containers of the stack, as \file{docker-compose.yml} defines them}\label{tab:ops:services}
\begin{tabularx}{\textwidth}{@{}L{35mm}L{33mm}Y@{}}
\toprule
\hd{Container} & \hd{Image} & \hd{What it does} \\
\midrule
\code{legalops-endpoint} & \code{caddy:2.11-alpine} & The only way in, on host port 80. Sends \code{/api/assistant/*} to the assistant service, \code{/api/*} to the API, \code{/dbgate/*} to the database tool after asking the API whether the caller is an administrator, and everything else to the console and the public site. Writes a JSON access log. \\
\code{legalops-exit} & \code{vuknf/caddy-l4:2.11} & The only way out: forwards TLS connections to a fixed list of host names (\cref{sec:ops:secrets}). \\
\code{legalops-vite} & built from \file{vite/} & The public site and the console: the Vite development server, or in production mode the built files served by Caddy. \\
\code{legalops-django} & built from \file{django/} & The API (Django). The worker's loop runs inside this container too. \\
\code{legalops-worker} & \code{docker:27-cli} & Nothing but one command, \texttt{docker exec legalops-django python3 manage.py run\_worker}: it starts the worker loop in the API's container and restarts it when it ends. \\
\code{legalops-assistant} & built from \file{assistant/} & The assistant service: the chat and every model call. \\
\code{legalops-postgres} & built from \file{postgres/} & PostgreSQL 17.11 with the pgvector extension 0.8.6 compiled in. \\
\code{legalops-dbgate} & \code{dbgate/dbgate:7.3.3-alpine} & A database administration tool for administrators, reached at \code{/dbgate/}. \\
\bottomrule
\end{tabularx}
\end{table}

\begin{figure}[H]
\centering
\begin{tikzpicture}[x=1cm, y=1cm]
  \node[ext]   (proxy)    at (0,1.2)     {outer proxy\\TLS};
  \node[svc]   (endpoint) at (2.75,1.2)  {endpoint\\port 80};
  \node[svc]   (vite)     at (6.1,2.45)  {vite\\site and console};
  \node[svc]   (asst)     at (6.1,1.2)   {assistant\\model calls};
  \node[svc]   (django)   at (6.1,-0.15) {django\\API + worker loop};
  \node[svc]   (worker)   at (6.1,-1.75) {worker\\Docker CLI only};
  \node[svc]   (exit)     at (9.75,1.2)  {exit\\allowlist};
  \node[store] (pg)       at (9.75,-0.15) {postgres};
  \node[svc]   (dbgate)   at (9.75,-1.75) {dbgate};
  \node[ext]   (gw)       at (12.75,1.75) {the gateway};
  \node[ext]   (src)      at (12.75,0.65) {sources of\\the acts};

  \draw[flow] (proxy) -- (endpoint);
  \draw[flow] (endpoint) -- (asst);
  \draw[flow] (endpoint.east) -- ++(0.55,0) |- (vite.west);
  \draw[flow] (endpoint.east) -- ++(0.55,0) |- (django.west);
  \draw[flow] (endpoint.south) -- (2.75,-2.75) -| node[lbl, pos=.25] {/dbgate/, administrators only} (dbgate.south);
  \draw[flow] (asst) -- (exit);
  \draw[flow] (django) -- (pg);
  \draw[flow] (dbgate) -- (pg);
  \draw[flow] ([yshift=2.5mm]django.east) -- ++(0.45,0) |- ([yshift=-2mm]exit.west);
  \draw[flow] (exit.east) -- ++(0.45,0) |- (gw.west);
  \draw[flow] (exit.east) -- ++(0.45,0) |- (src.west);
  \draw[call, <->] (django) -- node[lbl, right=1mm] {internal door} (asst);
  \draw[call] (worker) -- node[lbl, right=1mm] {docker exec} (django);

  \node[store] (d1) at (1.7,-4.9)   {\_DATA/postgres\\the database};
  \node[store] (d2) at (5.55,-4.9)  {\_DATA/django/uploads\\uploaded files};
  \node[data]  (envf) at (9.3,-4.9) {.env\\secrets};
  \begin{scope}[on background layer]
    \node[lane, fit=(endpoint)(vite)(django)(asst)(dbgate)(pg)(exit)(worker), inner ysep=9pt] (stack) {};
    \node[lane, fit=(d1)(d2)(envf), fill=codebg] (disk) {};
  \end{scope}
  \node[lanelabel, anchor=south west] at (stack.north west) {one Docker Compose project: eight containers};
  \node[lanelabel, anchor=south west] at (disk.north west) {on the host's disk, in the project folder};
\end{tikzpicture}
\caption{The stack on its host. Requests come in through the endpoint only, and connections to
the outside leave through the exit only. The worker container starts the worker loop inside
the API's container and has no network of its own. Everything the system remembers lies in
folders of the host that are mounted into the database and the API.}\label{fig:ops:host}
\end{figure}

\paragraph{Networks.} The compose file defines \code{isolated}, an internal bridge network on
172.28.10.0/24 that carries every container except the worker and has no route out, and joins
a second network, \code{external}, which it expects to exist already: the update script
creates it with \texttt{docker network create}, on the subnet 172.18.0.0/24, when the host has
no network of that name; a host that has one keeps it as it is, and the stack joins it. On the
live host the network was there before the stack, and other containers of the hosting are
attached to it too (seen on 16~November 2026). The endpoint and
the exit are on both; in development mode the console and the assistant service are on
\code{external} as well, so that they can install their packages.

\paragraph{State.} Everything the system remembers is in two folders that
\file{runUpdateThisStack.sh} creates and gives to user 1000: \file{_DATA/postgres}, the
database cluster, and \file{_DATA/django/uploads}, the files users and administrators upload.
The secrets are in \file{.env}, which Compose reads when it creates a container. None of the
three is in the repository.

\paragraph{What the host must offer.} According to \file{docs/markdown/en/deployment.md}:
\begin{itemize}
  \item Linux with Docker Engine (tested on 27.3.1) and both front ends of Compose: the update
    script calls \code{docker-compose}, the other scripts \texttt{docker compose}. The update
    script also uses \code{sudo} and \code{openssl}.
  \item Four processor cores and 8~GB of memory are comfortable. No container has a memory
    limit; the peaks are the sandboxed children that extract text from uploaded files (1.5~GB
    of address space each, \code{CHILD_MEMORY_MB} in \file{common/sandbox.py}) and the
    rendering of a report as PDF.
  \item About 3~GB of disk for the images, plus the data.
  \item Host port 80 free, and the two subnets above unused.
  \item Outbound HTTPS to the gateway (\code{ai.knf.vu.lt}) and to the three sources of the
    acts (\code{publications.europa.eu}, \code{www.e-tar.lt}, \code{www.lb.lt}); while images
    are built, also to the package registries.
  \item A reverse proxy in front that terminates TLS, because the endpoint speaks plain HTTP.
    It must pass on the \code{Host} header as the browser sent it, since the endpoint compares
    a request's \code{Origin} with it, and put the visitor's address, and only that, into
    \code{X-Forwarded-For}, which the endpoint copies into \code{X-Real-IP} (what the live
    proxy sends is in \cref{ch:security}). It must not buffer \code{text/event-stream} (the
    chat streams its answer), must wait at
    least 330~s for an answer (a report export is one long request, and a prompt test run may
    wait up to fifteen minutes: a proxy that cuts it earlier shows an error in the prompt
    editor while the call still completes and is recorded),
    must let through request bodies of 62~MiB (the API accepts an uploaded act of up to
    60~MiB) and must leave the endpoint's two security headers in place.
\end{itemize}

\paragraph{Capacity.} In production mode the API runs under gunicorn with
\code{GUNICORN_WORKERS} processes of \code{GUNICORN_THREADS} threads each (3 and 8 by
default, \file{django/gunicorn.conf.py}), so 24 requests are served at once. One worker loop
runs one job at a time; inside a job up to four model calls are in flight
(\cref{app:config}). The queue is claimed with \texttt{SELECT … FOR UPDATE SKIP LOCKED}
(\code{claim} in \file{jobs/runner.py}), so a second loop would be safe; it needs its own
\code{WORKER_PID_FILE}, because a loop that finds another loop's PID file takes over from it
instead of running beside it. \file{docs/markdown/en/operations.md} names the real ceiling: the gateway's
budget, not the host.


\section{A fresh installation}\label{sec:ops:install}

An installation is four things in this order: the two local files, one run of the update
script, the first administrator, and one re-embedding of the corpus. The legal corpus itself
is not fetched and not read by a model again: it arrives with the database's bootstrap dump.

\subsection{The local files and the update script}

\begin{codeblock}
git clone <repository> /apps/aplikacijos/comp-assistant
cd /apps/aplikacijos/comp-assistant
cp docker-compose.sample.yml docker-compose.yml
touch .env && chmod 600 .env      # then add the line  AI_GATEWAY_KEY=<the gateway key>
./runUpdateThisStack.sh
\end{codeblock}

The repository versions \file{docker-compose.sample.yml}; \file{docker-compose.yml} is the
local copy an installation runs and is ignored by git, as \file{.env} is.

\begin{notebox}{Two lines of the compose file an installation depends on}
The copy is the installation's own: no script makes it, and \texttt{git pull} never changes it,
so a later change of the sample reaches it only when someone carries it over. Two lines of the
sample matter most to a new installation: the read-only mount of
\file{postgres/bootstrap.sql.gz} into the database container's
\file{/docker-entrypoint-initdb.d}, and the variable \code{DBGATE_AUTH_HEADER} on the
endpoint. Without the first, a new installation starts with an empty database: the update
script's \code{migrate} then creates the schema, and the operator has to load the prompts
(\code{seed_prompts}, \cref{sec:ops:prompts}) and the catalogue (\code{seed_corpus},
\cref{sec:ops:corpus}) and pay for every reading of the corpus again. Without the second, the
endpoint cannot sign in to the database tool. The repository holds no sample of \file{.env}:
the file starts as the one line with the gateway key, and the update script appends the
secrets it mints.
\end{notebox}

\file{runUpdateThisStack.sh} does the same six things on every run, the first one included:

\begin{enumerate}
  \item Creates \file{_DATA/django/uploads} and \file{_DATA/postgres} and gives \file{_DATA}
    to user 1000.
  \item Mints \code{DJANGO_SECRET_KEY} (32 random bytes as hexadecimal, from
    \texttt{openssl rand}) and appends it to \file{.env} --- only when \file{.env} has no line
    that begins with that name. A line that is there, even an empty one, is left alone.
  \item Mints \code{ASSISTANT_INTERNAL_SECRET} the same way. It is the shared secret of the
    internal door between the API and the assistant service; with an empty value both sides
    refuse every internal call.
  \item Mints \code{DBGATE_PASSWORD}, the password of the database tool's own login, and
    \code{DBGATE_AUTH_HEADER}, the same login encoded for an HTTP header, which the endpoint
    adds to every request it passes to the tool.
  \item Creates the Docker network \code{external} (an error whenever a network of that name
    exists --- on every later run, or from the start on a host where something else made it
    --- and harmless), takes the stack down and starts it again from freshly built images:
    \texttt{docker-compose down}, then \texttt{docker-compose up -d -{}-build}.
  \item Waits five seconds and runs \texttt{manage.py migrate -{}-noinput} in the API's
    container.
\end{enumerate}

The script never writes \code{AI_GATEWAY_KEY}: the key of the faculty's gateway is the one
value a person supplies. Without it the stack starts, but the assistant service reports its
model as not configured, the chat answers 503 and every job step that needs a model fails.

\subsection{The bootstrap dump}\label{sec:ops:bootstrap}

When the database container starts on an empty \file{_DATA/postgres}, the official PostgreSQL
entrypoint initialises a cluster and loads every file it finds in
\file{/docker-entrypoint-initdb.d} --- here one file, \file{postgres/bootstrap.sql.gz}, about
11~MB of gzipped plain SQL. A cluster that already exists is never touched by it. The file is
written by \file{scripts/db_bootstrap_dump.sh} from a running installation
(\cref{fig:ops:bootstrap}).

\begin{figure}[H]
\centering
\begin{tikzpicture}[node distance=7mm and 6mm]
  \node[store] (live) {the live\\database};
  \node[svc, right=of live] (copy) {a throwaway\\postgres\\no network};
  \node[chk, right=of copy] (cut) {the cut: keep\\eight tables,\\clear the rest};
  \node[data, right=of cut] (file) {bootstrap\\.sql.gz};
  \node[store, below=13mm of live] (new) {an empty\\data folder};
  \node[svc, right=of new] (load) {the entrypoint\\loads the file};
  \node[man, right=of load] (admin) {the first\\administrator};
  \node[ai, right=of admin] (embed) {one\\re-embedding};
  \draw[flow] (live) -- node[lbl, above] {read} (copy);
  \draw[flow] (copy) -- (cut);
  \draw[flow] (cut) -- (file);
  \draw[flow] (new) -- (load);
  \draw[flow] (load) -- (admin);
  \draw[flow] (admin) -- (embed);
  \draw[opt] (file.south) |- ([yshift=6mm]load.north) -- (load.north);
  \node[note, anchor=west] at ([xshift=3mm]file.east) {made on the\\running system};
  \node[note, anchor=west] at ([xshift=3mm]embed.east) {on the new\\installation};
\end{tikzpicture}
\actorlegend
\caption{How a new installation gets its corpus. Above: \file{scripts/db_bootstrap_dump.sh}
copies the live database into a throwaway container, cuts it down there and dumps the
result; the live database is only read. Below: a new installation loads that file once,
creates its first administrator, who then starts the one paid step --- the vectors, which the
file leaves out.}\label{fig:ops:bootstrap}
\end{figure}

\paragraph{What the file holds.} The script keeps an allow-list of eight tables (the array
\code{keep} in its SQL): \code{django_migrations}, \code{legal_documents},
\code{legal_document_versions}, \code{legal_units}, \code{legal_chunks}, \code{obligations},
\code{prompts} and \code{prompt_versions}. A new installation therefore starts with
\begin{itemize}
  \item the schema, with every migration recorded as applied;
  \item the legal corpus: the acts of the catalogue, the current version of each act that has
    one, its provisions and its chunks, and whatever was read into the version rows (the scope
    and sanction readings of \cref{ch:obligations} are columns of
    \code{legal_document_versions});
  \item the register of obligations with the judge's verdicts and the review statuses ---
    5,960 obligations of 68 acts in the file shipped on 19~November 2026;
  \item the prompts with their versions, as they stood on the day of the dump: the file of
    20~November 2026 holds 91 versions, the chat's prompt in its seventeenth.
\end{itemize}

\paragraph{What it leaves out.} Every other table is emptied: accounts, sessions, companies
and everything under them, reports, chats, work products, jobs, the audit log, the record of
model calls, the price cache. Because the list names what stays, a table added to the system
later ships empty until someone adds it to that list --- nothing leaks by default. Inside
the kept tables the script removes what belongs to people or to one installation: a company's
own requirements (obligations with a company), the links from kept rows to accounts
(\code{created_by}, \code{reviewed_by}), superseded versions of acts with their units and
chunks, and every embedding. The embeddings go because of their size --- as SQL text they
were three quarters of a 156~MB file, by \file{docs/markdown/en/operations.md} --- and
\code{embed_model} is cleared with them, so the console shows which versions still wait for
their vectors. The raw source files of the acts are not included either; nothing reads them
once a version is ready. The cut runs in one transaction with all foreign keys deferred, so a
kept row that still points at an emptied table fails the dump on the running system and not
the installation that would load it.

The load takes about ten seconds by \file{docs/markdown/en/deployment.md},
longer than the update script's five-second wait. On a first run the script's \code{migrate}
can therefore fail with \enquote{connection refused}; the dump is already migrated, so the
failure is harmless then.

\subsection{The first administrator}\label{sec:ops:admin}

Nobody can register themselves, and every route that creates an account needs an
administrator. The way into an empty installation is the management command
\code{create_admin} (\file{users/management/commands/create_admin.py}):

\begin{codeblock}
read -rsp 'Admin password: ' ADMIN_PASSWORD; echo; export ADMIN_PASSWORD
docker compose exec -T -e ADMIN_PASSWORD legalops-django \
    python3 manage.py create_admin --email <address> --password-env ADMIN_PASSWORD
unset ADMIN_PASSWORD
\end{codeblock}

\file{runCreateAdmin.sh} in the repository's root asks for the address and the password and
runs exactly this. The password travels in an environment variable named by \code{--password-env}, never on the
command line, where a process listing would show it. Unlike a password set in the console,
it has no minimum length: whoever runs a command on the server chooses it, and the command
refuses only an empty one and one past bcrypt's 72 bytes (\code{PASSWORD_MAX_BYTES}). The
command does nothing to an account that exists. With \code{--reset} it makes the account an
active administrator again and, when a password is given, sets it and ends every session of
the account; that is the recovery from a lock-out, and it is written to the audit log
(\code{users.reset_password}, or \code{users.update} when no password is given) with no actor.

\subsection{The one re-embedding}

The administrator signs in, opens the corpus page (\ui{Teisės aktų bazė}, \enquote{legal
acts}) and chooses \ui{Perskaičiuoti visus vektorius} (\enquote{re-embed everything}) in its
menu. That queues one job, \code{corpus.reembed}, through
\route{POST}{/api/corpus/reembed}. The job (\code{reembed} in \file{corpus/ingest.py}) sends
the chunks of every act's current version to the embedding model in batches of 64
(\code{EMBED_BATCH}), then the Lithuanian and the English text of every obligation, then the
chunks of whatever company documents exist --- none on a new installation. It is charged to
nobody. On 13~November 2026 the corpus held 4,059 chunks and 5,960 obligations;
\file{docs/markdown/en/deployment.md} puts the job, for about 4,100 chunks and 11,900 obligation
texts, at several minutes and under one dollar, which is its estimate and not a dated
measurement. \Cref{ch:models} describes the embedding call.

Until the job has run, three things are visibly wrong, and an operator should know them as
the signs of a skipped step and not as faults: the search of the corpus finds keyword matches
only, a report on a library too large to be sent whole finds no evidence for its obligations
(a library that fits is read whole and needs no vectors; \cref{ch:report}), and the table of an
act's versions
shows a dash where the embedding model's name belongs.

The same job is the answer whenever \code{AI_EMBED_MODEL} changes. Vectors of two models
cannot be compared, so the job renews all three sets together; the variable is set in two
places (the API's copy is written in \file{docker-compose.yml} itself, the assistant
service's comes from \file{.env}, with the same model as its default) and both must name the same model.

\subsection{Checking the result}

\begin{codeblock}
docker compose ps                       # every container Up
curl -s http://127.0.0.1/api/health     # {"status": "ok", "service": "legalops-django", ...}
\end{codeblock}

The bootstrap is
as current as its last dump: when the repository's prompt defaults or its catalogue are newer
than the file, \code{seed_prompts} and \code{seed_corpus} add what is missing, and the weekly
refresh brings the acts up to date.


\section{Development and production modes}\label{sec:ops:modes}

\file{docker-compose.yml} holds both modes in one file. Lines that belong to one mode carry a
trailing comment, \texttt{\# Dev} or \texttt{\# Prod}; the file as shipped has the development
lines active. To switch to production an operator comments out every line tagged
\texttt{\# Dev}, removes the comment sign from every line tagged \texttt{\# Prod}, and runs
\file{runUpdateThisStack.sh}. \Cref{tab:ops:modes} lists what changes.

\begin{table}[H]
\caption{What the two modes of the compose file run}\label{tab:ops:modes}
\begin{tabularx}{\textwidth}{@{}L{27mm}YY@{}}
\toprule
\hd{Service} & \hd{Development (as shipped)} & \hd{Production} \\
\midrule
Console and public site & \file{vite/Dockerfile.dev}: \texttt{npm install} and the Vite development server on the bind-mounted \file{vite/app}; also on the \code{external} network. & \file{vite/Dockerfile}: the page and its scripts are built in the image (\texttt{npm ci}, the translation guard, \texttt{vite build}) and served by Caddy; no source mount, isolated network only. \\
API & \code{DJANGO_RUNSERVER} is \code{True}: Django's \code{runserver} on the bind-mounted \file{django}, after waiting for the database. & The variable is unset: gunicorn from the code baked into the image, tuned by \file{django/gunicorn.conf.py}. \\
Assistant service & \file{assistant/Dockerfile.dev}: \texttt{npm install} and \code{nodemon} on the bind-mounted \file{assistant/app}; also on the \code{external} network. & \file{assistant/Dockerfile}: \texttt{npm ci -{}-omit=dev} and \texttt{node server.js}; a 64~MB temporary filesystem, isolated network only. \\
\bottomrule
\end{tabularx}
\end{table}

Two consequences matter to an operator. In development mode a saved source file reloads the
console, the assistant service and the API at once, but not the worker loop:
\texttt{docker compose restart legalops-worker} starts a loop on the new code. And in
development mode the assistant service has a second network with a route out, so its traffic
is not held to the exit's allowlist; production mode ends that. In production a code change
reaches the containers only through the update script. \Cref{ch:security} weighs what the
development mode costs in protection.


\section{Updating}\label{sec:ops:update}

\begin{codeblock}
git pull
./runUpdateThisStack.sh
docker compose exec legalops-django python3 manage.py seed_prompts
\end{codeblock}

The script is the one of the installation: it takes the whole stack down, rebuilds the
images (cached layers make a change of code alone a matter of seconds), starts everything
and migrates the database. The site is unreachable from the \code{down} until the containers
are back; building first with \texttt{docker compose build} keeps the build out of that window.
The third line is not optional when the update changed a file under
\file{prompts/defaults/}: nothing runs \code{seed_prompts} on its own
(\cref{sec:ops:prompts}). The test suites are not run by the script either
(\cref{ch:testing}).

Docker gives a stopping container ten seconds, since the compose file sets no
\code{stop_grace_period}. What that means for work in flight:

\begin{itemize}
  \item \textbf{A chat turn.} The assistant service stops taking connections at once, refuses
    a request on a connection kept open from before with 503 \code{SHUTTING_DOWN}, and lets
    running turns finish --- its own patience,
    \code{SHUTDOWN_GRACE_MS}, is 190~s --- but Docker ends the process after ten seconds, and
    a turn still streaming then loses its answer.
  \item \textbf{A request to the API.} Requests still running after ten seconds are cut.
  \item \textbf{A job.} The worker loop is a second process in the API's container and dies
    with it. The job it was running stays marked as running until the next loop finds its
    heartbeat older than \code{STALE_S} (120~s, \file{jobs/runner.py}); it is then queued
    again if its kind allows another attempt, and marked failed with \enquote{The worker
    running this job stopped unexpectedly} if not.
\end{itemize}

\Cref{tab:ops:jobs} shows which kinds get a second attempt. Most of the paid kinds do not, on
purpose: a job that has spent money is not silently run twice. An operator therefore looks at
the jobs page of the console (\cref{ch:console}) before an update, and does not update while
a report or a reading of the corpus is running. A report whose job ends this way is marked
failed by the kind's \code{on_abandon} hook; what a failed report keeps is in
\cref{ch:report}.

\begin{table}[H]
\caption{The kinds of job, how many attempts each is registered with (\code{max_attempts})
and what starts it}\label{tab:ops:jobs}
\begin{tabularx}{\textwidth}{@{}L{48mm}C{16mm}Y@{}}
\toprule
\hd{Kind} & \hd{Attempts} & \hd{Started by} \\
\midrule
\code{report.analyze} & 1 & A user who starts a compliance report, and again when its scope is approved. \\
\code{programme.generate} & 1 & A user who asks for a work product. \\
\code{documents.index} & 3 & Every upload of a document that has text. \\
\code{corpus.ingest} & 2 & \code{seed_corpus}, a refresh, an act added on the corpus page, an uploaded version of an act. \\
\code{corpus.refresh_all} & 1 & The weekly schedule; the corpus page. \\
\code{corpus.reembed} & 1 & The corpus page. \\
\code{corpus.extract_obligations} & 1 & The command of that name; an act's page. \\
\code{corpus.verify_obligations} & 1 & The command of that name; an act's page. \\
\code{corpus.read_sanctions} & 1 & The command of that name. \\
\code{ops.maintenance} & 1 & The daily schedule. \\
\bottomrule
\end{tabularx}
\end{table}

A worker restarted on its own behaves better than one killed with its container. A new loop
reads the PID file of the loop still running and sends it \code{SIGTERM}
(\code{take_over} in \file{jobs/management/commands/run_worker.py}); the old loop stops
claiming, its handler stops at its next cancellation check, and its job goes back to the
queue \emph{without} spending an attempt.

The base images are pinned by tag, not all to the same depth: \code{caddy:2.11-alpine},
\code{vuknf/caddy-l4:2.11} and \code{python:3.13-slim} follow a minor line and
\code{docker:27-cli} a major line, while \code{node:24.16.0-alpine} (\code{node:24.16.0} in
the two development images), \code{postgres:17.11-alpine3.24} and
\code{dbgate/dbgate:7.3.3-alpine} name an exact release, so that a newer release of those
arrives only with a changed tag.
\file{docs/markdown/en/operations.md} asks for one rebuild a month with \texttt{docker compose build -{}-pull}
before the update script, to take in what was published anew under the tags of the images
built here; an image Compose only pulls --- the endpoint's, the exit's, the worker's, the
database tool's --- is refreshed with \texttt{docker compose pull}. A new minor version of
PostgreSQL 17 is a changed tag in \file{postgres/Dockerfile} and a rebuild; a new major
version needs the database dumped and loaded into an empty cluster of the new version, which
nothing in the repository automates.


\section{Operating the corpus}\label{sec:ops:corpus}

The corpus needs four kinds of attention from an operator: bringing in acts the catalogue
gained, keeping the acts current, giving the system a file where it cannot fetch one, and
starting the three readings by a model that cost money, which nothing starts over the whole
corpus by itself (for one act an administrator can start two of them from its page, and a report
has an act it needs read).
What ingestion and the readings do inside is the subject of \cref{ch:corpus,ch:obligations};
this section says when to start them, what they queue and what they cost.

\subsection{Seeding, refreshing, uploading}

\paragraph{Seeding the catalogue.} \code{seed_corpus}
(\file{corpus/management/commands/seed_corpus.py}) reads the catalogue,
\file{corpus/catalog.yaml} --- 78 acts, 68 of them enabled --- validates each entry, and
creates the act rows that do not exist yet. For each new, enabled act with a source that can
be fetched it queues one \code{corpus.ingest} job. Acts that exist are left alone: after the
first seed the database is the truth, not the file.

\begin{codeblock}
docker compose exec legalops-django python3 manage.py seed_corpus                   # every missing act
docker compose exec legalops-django python3 manage.py seed_corpus --only mica,dora  # only these slugs
docker compose exec legalops-django python3 manage.py seed_corpus --no-ingest       # rows only
\end{codeblock}

An installation that started from the bootstrap dump needs the command only when the
catalogue is newer than the dump. One caution follows from how it works: it also creates
again an act that an administrator deleted on purpose; \code{--only} limits it. Ingestion
spends a little on embeddings, charged to nobody: bringing in 24 new catalogue acts on a copy
of the database on 12~November 2026 --- 22 of them fetched, and the two guidelines that have
to be uploaded (below) --- took 89~s and \$0.08.

\paragraph{Refreshing.} A refresh fetches an act from its source again; a new version is
stored only of a text that changed (\cref{sec:corpus:unchanged}), so an unchanged act costs
one request and no model call. An EU act whose version index does not answer is left as it
is, with the failure as its last error, and a new version whose text the act has had before
takes over the obligations and readings of that earlier version (\cref{ch:corpus}). The worker's schedule queues \code{corpus.refresh_all} once a week
(\cref{sec:ops:schedule}); that job queues one \code{corpus.ingest} for every enabled act
whose source is not \code{upload}. An administrator starts the same from the corpus page
(\ui{Tikrinti visus šaltinius}, \enquote{check all sources};
\route{POST}{/api/corpus/refresh-all}) or for one act from its page
(\route{POST}{/api/corpus/documents/{id}/refresh}). The fetchers pace themselves to one
request a second per host (\code{MIN_INTERVAL_S} in \file{corpus/fetchers.py}) and use HTTPS
only.

\paragraph{An act that cannot be fetched.} Two enabled acts of the catalogue --- the European
Banking Authority's guidelines on money-laundering risk factors and their amendment for
crypto-asset service providers --- have the source \code{upload}: the authority's site is not
on the exit's allowlist, so nothing is fetched for them. \code{seed_corpus} creates their rows
and prints them as waiting for a file. An administrator opens each act on the corpus page,
downloads the Lithuanian PDF from the act's source link and uploads it with \ui{Įkelti naują
redakciją} (\enquote{upload a new version}); the upload becomes a pending version and its
ingestion is queued. The same button serves any act whose source fails, and it is how a
licensed standard enters the corpus (\cref{ch:corpus}). A file may be up to 60~MiB
(\code{MAX_UPLOAD_BYTES} in \file{corpus/api/views.py}).

\subsection{The three paid readings}\label{sec:ops:readings}

Three management commands hand the corpus to a model. Each only \emph{queues} jobs and
returns at once; the worker does the reading, and the jobs page shows the progress and the
errors. \Cref{tab:ops:readings} sets them side by side, and \cref{fig:ops:readings} shows the
order in which an operator uses them.

\begin{figure}[H]
\centering
\begin{tikzpicture}[node distance=6mm and 5mm]
  \node[svc] (ingest) {acts ingested\\\code{seed_corpus},\\refresh, upload};
  \node[ai, right=of ingest] (extract) {\code{extract_}\\\code{obligations}};
  \node[ai, right=of extract] (verify) {\code{verify_}\\\code{obligations}};
  \node[man, right=of verify] (review) {review on the\\obligations page};
  \node[ai, right=of review] (sanctions) {\code{read_}\\\code{sanctions}};
  \node[svc, below=14mm of sanctions] (dump) {\code{db_bootstrap_}\\\code{dump.sh}};
  \draw[flow] (ingest) -- (extract);
  \draw[flow] (extract) -- (verify);
  \draw[flow] (verify) -- (review);
  \draw[flow] (review) -- (sanctions);
  \draw[opt] (sanctions) -- (dump);
  \node[note, below=1mm of extract, align=center] {a model reads every\\provision for duties};
  \node[note, below=1mm of verify, align=center] {a judge compares each\\duty with its provisions};
  \node[note, below=1mm of review, align=center] {a person approves,\\corrects or rejects};
  \node[note, left=3mm of dump, align=right] {so that new installations\\inherit the readings};
\end{tikzpicture}
\actorlegend
\caption{The paid readings of the corpus in the order an operator starts them. The reading of
sanctions does not depend on the other two and can run at any time; it stands last here
because it is the cheapest and the least urgent.}\label{fig:ops:readings}
\end{figure}

\begin{table}[H]
\caption{The three paid readings an operator starts. Every job is queued at priority $-1$ and
charged to nobody when started from the command line. The costs are single dated runs on the
live system with \model{gpt-6.1-sol}.}\label{tab:ops:readings}
\begin{tabularx}{\textwidth}{@{}L{25mm}YYY@{}}
\toprule
 & \hd{\code{extract_obligations}} & \hd{\code{verify_obligations}} & \hd{\code{read_sanctions}} \\
\midrule
\hd{Prompt} & \prompt{corpus.extract_obligations}; for a standard \prompt{corpus.extract_requirements} & \prompt{corpus.verify_obligations} & \prompt{corpus.sanctions} \\
\hd{Queues} & One job per act. & One job per act. & One job for all the acts. \\
\hd{Without arguments} & Every enabled act that is \emph{behind}: never read, or read before the prompt version now active was written. & Every enabled act with obligations no judge has seen. & Every enabled act never read for sanctions; standards are left out. \\
\hd{With slugs} & Those acts, whatever their state. & Those acts, as far as they hold unjudged obligations; with \code{--again} all their obligations. & Those acts, as far as they were never read; with \code{--again} all of them. \\
\hd{Order} & Smallest act first. & Fewest obligations first. & By jurisdiction, then slug. \\
\hd{Measured} & 10~Nov 2026: 51 acts, 461 calls, \$23.61, two and a half hours. 13~Nov: 17 acts, 36 calls, \$4.77. & 13~Nov 2026: 5,960 obligations, 369 calls, \$9.80. & 13~Nov 2026: 29 calls, \$0.71, 262 tiers in 27 acts. 19~Nov, read anew on a copy: 29 calls, \$0.61, 309 tiers with version~2 of the prompt; 29 calls, \$0.63, 288 tiers with version~3. 20~Nov, with version~4: 29 calls, \$0.61, 302 tiers. \\
\bottomrule
\end{tabularx}
\end{table}

\begin{codebox}{What the three commands check before anything is spent}
All three are written the same way (\file{corpus/management/commands/}). Each looks up the
active version of its prompt and stops with \enquote{run seed\_prompts first} when there is
none. Each prints the prompt version and the model the acts will be read with, so that a
prompt that is not the expected one shows before a cent is spent. Each refuses to queue a
second job for an act (or, for the sanctions, a second job at all) while one is queued or
running. And each takes by default only what is still missing, so that running the command
again after a failure --- a spent gateway budget, an outage --- finishes the work and pays for
nothing twice.
\end{codebox}

\paragraph{Priorities.} The worker claims the queued job with the highest priority first and
the oldest among equals (\code{claim} in \file{jobs/runner.py}). Work that users start has
priority 0 (the indexing of an uploaded document 5), the three commands use $-1$ (\code{PRIORITY}) and the schedule $-5$. A report
queued in the middle of a re-reading of the whole corpus therefore runs before the next act
is read, though not before the act in hand is finished: one loop runs one job at a time.

\paragraph{Extraction --- when the acts must be read again.} The command is \code{extract_obligations}. The obligations
in the register are the work of one prompt version on one model. After the extraction prompt
or its model changes, every act is \emph{behind} (\code{is_behind} in
\file{corpus/obligations.py} compares the time an act's version was read with the time the
active prompt version was written), and the command queues them all. Acts new to the corpus
are behind as well, although the first report that needs such an act reads it by itself; the
command reads them ahead of any report.

\begin{codeblock}
docker compose exec legalops-django python3 manage.py extract_obligations             # every act that is behind
docker compose exec legalops-django python3 manage.py extract_obligations mica dora   # these acts, whatever their state
\end{codeblock}

A job of the bulk run carries the flag \code{only_behind}: if a report has read the act by the
time the job's turn comes, the job ends as skipped and nothing is paid twice. An act keeps
its present proposed obligations until its new extraction has succeeded as a whole, and
obligations an administrator has approved or rejected are never touched.

\begin{humanbox}{Review before re-reading}
The action register (\cref{ch:register}) knows an obligation by its identifier, and an
obligation keeps its identifier across a re-reading of its act only once a person has
approved it; a proposed one is replaced. An operator therefore has the register reviewed
before re-extracting an act whose obligations stand under open actions --- otherwise those
actions can only be closed by hand. This is a decision about order that the system leaves to
its operator.
\end{humanbox}

\paragraph{The judge --- before a review of the register.} The command is \code{verify_obligations}. The judge, a
second model call, compares each obligation with the full text of the provisions it cites and
marks it supported, partly supported or unsupported, with a note (\cref{ch:obligations}). It
approves and rejects nothing: its verdicts only put the doubtful entries first on the
obligations page, so that a person's review starts where it is needed. The command is run
after an extraction and before that review, and again with \code{--again} after the judge's
prompt or model changed.

\begin{codeblock}
docker compose exec legalops-django python3 manage.py verify_obligations             # every act with unjudged obligations
docker compose exec legalops-django python3 manage.py verify_obligations --again mica  # judge these anew
\end{codeblock}

\paragraph{Sanctions --- what a breach can cost.} The command is \code{read_sanctions}. A report reads the sanction
provisions of the acts it assesses by itself. The command is for the rest of the corpus: a
national law that lays down the fines for an EU regulation is found by a finding under that
regulation only once the national law has been read. It is run once after acts were added,
and with \code{--again} after the prompt changed. The one job reads the acts one after
another and stops at the first that cannot be read; the acts read so far keep their result,
and the next run takes up the rest.

\begin{codeblock}
docker compose exec legalops-django python3 manage.py read_sanctions           # every enabled act never read
docker compose exec legalops-django python3 manage.py read_sanctions --again   # read all anew
\end{codeblock}

\begin{lessonbox}{Measured: what a whole corpus costs to read}
On 10~November 2026 all 51 acts then enabled were read again for obligations with
\model{gpt-6.1-sol}: 461 model calls, \$23.61, from 12:42 to 15:12~UTC, no failure, 4,892
obligations. On 13~November the 17 acts added that week cost 36 calls and \$4.77 and gave
1,068 more. The judge then read all 5,960 obligations in 369 calls for \$9.80, and the
sanction provisions of the whole corpus took 29 calls and \$0.71. Together the three readings
of a corpus of 68 acts cost about \$39 at those prices, \$28.38 of it the extraction in its
497 calls --- more than the gateway's team budget
held on the day it ran out (\cref{sec:ops:budgets}).
\end{lessonbox}

\subsection{Refreshing the bootstrap dump}

After a refresh of the corpus, an extraction or a review that new installations should
inherit, the operator writes the dump again:

\begin{codeblock}
scripts/db_bootstrap_dump.sh
\end{codeblock}

It does what \cref{sec:ops:bootstrap} describes; by
\file{docs/markdown/en/operations.md} it took about 40~s when the file was 9~MB. The live database is only read. The new file is committed with the
code, which is how the readings paid for once reach every later installation.


\section{Operating the prompts}\label{sec:ops:prompts}

Prompts are rows in the database, edited in the console (\cref{ch:prompts}); the repository
ships their defaults as YAML files in \file{prompts/defaults/}, 23 of them. Two commands
connect the files, the database and the models.

\subsection{Loading the shipped prompts after an update}

\code{seed_prompts} calls \code{seed_defaults} in \file{prompts/store.py}, which reads every
default file and is safe to run any number of times:

\begin{itemize}
  \item a prompt that does not exist is created with the default as its active version;
  \item a prompt's name, description, variables and output schema always follow the file,
    because the code reads them;
  \item a default not seen before (recognised by a hash of its text, model, temperature and
    output limit, \code{default_hash}) becomes a new prompt version. It is made active at once while the prompt still runs a shipped default. Where
    an administrator has activated a text of their own, the new default waits beside it as a
    version to compare and activate by hand;
  \item where an administrator kept the default's text but chose another model or other
    settings, the new text is activated on those settings.
\end{itemize}

The command prints which prompts were seeded, which were upgraded and which wait. It is run
by hand after an update, and by nothing else. Forgetting it after an
update leaves the existing prompts on their previous texts without any error. A prompt the
update added does not exist until the command has run, and neither does any prompt of a
database that did not start from the bootstrap dump: a job step that needs one fails with
\enquote{Prompt '…' has no active version}.

One side effect is worth knowing before the command is run on a large corpus: a new active
version of the extraction prompt makes every act \emph{behind}, so the next
\code{extract_obligations} without arguments offers to read the whole corpus again
(\cref{sec:ops:readings}).

\subsection{Measuring before a prompt or a model changes}

The test suites of \cref{ch:testing} run against a fake model and prove the plumbing. What a
change of a prompt or of its model does to the model's answers is measured by
\code{eval_prompts} (\file{prompts/management/commands/eval_prompts.py}), which puts labelled
cases to the real model through the same functions the product calls (the labels are an AI
model's, not yet confirmed by a person). \Cref{ch:control}
describes the suites and their cases; this section is the operator's side.

\begin{codeblock}
docker compose exec legalops-django python3 manage.py eval_prompts            # every suite
docker compose exec legalops-django python3 manage.py eval_prompts verify     # one suite
docker compose exec legalops-django python3 manage.py eval_prompts --record   # accept the scores as the baseline
\end{codeblock}

\begin{table}[H]
\caption{The six suites of \code{eval_prompts} and the baseline recorded in
\file{prompts/evals/baseline.json} on 12~November 2026 (the sanctions suite again on
20~November, for version~4 of its prompt and with sixteen cases), all with
\model{gpt-6.1-sol}}\label{tab:ops:suites}
\begin{tabularx}{\textwidth}{@{}L{22mm}L{47mm}YL{22mm}@{}}
\toprule
\hd{Suite} & \hd{Prompt it measures} & \hd{What a case asks} & \hd{Baseline} \\
\midrule
\code{verify} & \prompt{corpus.verify_obligations} & Does this provision carry this obligation? & 14 of 14 \\
\code{scope} & \prompt{corpus.scope} & Whom does this act's own text cover? & 9 of 9 \\
\code{citations} & \prompt{corpus.check_citations} & Does this article regulate what a document says it does? & 19 of 20 \\
\code{gap} & \prompt{report.gap_analysis} & Do the sample company's documents meet this obligation? & 67 of 71, slack 2 \\
\code{applicability} & \prompt{report.applicability} & Does this act apply to this company? & 19 of 19 \\
\code{sanctions} & \prompt{corpus.sanctions} & What can a breach of this article cost? & 16 of 16 \\
\bottomrule
\end{tabularx}
\end{table}

For each suite the command prints every case the model got wrong, with what was expected and
what it answered, then the score beside the baseline's, and at the end what the model calls
since the run began cost (the sum of the recorded calls, so calls of other work running at
the same time are counted in). It ends with an error when a suite scores below its baseline.
A suite whose answers vary a little from run to run carries a \code{slack} in its baseline
entry, the number of cases a run may lose before it fails: 2 for the gap suite. With
\code{--record} the scores are written as the new baseline together with the prompt version,
the model and the date that earned them, and the command does not fail. The version is named
by its number in the database of the installation that recorded it, and that is by design: the
same wording may carry another number on another installation (\cref{ch:control}).

A full run of the six suites cost about a dollar when it was measured on 12~November 2026,
nearly half of it, \$0.42, the gap suite (\file{docs/markdown/en/operations.md}).
The \code{scope}, \code{citations}, \code{applicability} and \code{sanctions} suites read
this installation's corpus, so they need the acts their cases name; a case whose act is
missing counts as failed and says why. The gap suite brings its own company: the 26 policies
and the profile of \demo{} are kept beside the cases in \file{prompts/evals/data/1009-ltd},
so it judges the same documents on every installation.

\begin{humanbox}{Recording a baseline is a person's decision}
The procedure is: run the suites, change the prompt or the model, run them again, read every
miss. Only then does a person decide whether the new scores are the standard to hold future
changes to, and record them. A miss can also mean that the label is wrong; a label is
overruled by editing its case file, which is a change in the repository and visible in its
history.
\end{humanbox}


\section{Budgets and quotas}\label{sec:ops:budgets}

Four limits stand between a user and the gateway's bill or the server's disk, from the outermost
in.
\Cref{ch:models} explains how a call is priced and recorded; here are the limits an operator
sets and the symptoms of each.

\begin{table}[H]
\caption{The limits on spending and storage, and what happens at each}\label{tab:ops:limits}
\begin{tabularx}{\textwidth}{@{}L{30mm}L{36mm}Y@{}}
\toprule
\hd{Limit} & \hd{Set where} & \hd{When it is reached} \\
\midrule
The gateway's team budget & At the gateway, for the key; nothing in this system changes it. & The gateway answers HTTP 402 to every call. The chat shows \code{AI_BUDGET_EXHAUSTED}; job steps fail with the same code and are not retried. \\
The daily budget of a user & \code{LLM_DAILY_BUDGET_USD} (default 10.0; 0 means unlimited), or per account on the users page (\ui{Naudotojai}). & A chat turn is refused with 429 \code{BUDGET_EXCEEDED}; so are a prompt test run, and the start of a report or of a work product when less than \$0.25 is left. \\
The upload quota of a user & \code{UPLOAD_QUOTA_FILES} (1000) and \code{UPLOAD_QUOTA_BYTES} (2~GiB). & An upload is refused with 413 \code{QUOTA_EXCEEDED}. A deleted document frees its share. \\
The size of one file & \code{MAX_FILE_BYTES}, 50~MiB, for a document or a chat attachment; \code{MAX_UPLOAD_BYTES}, 60~MiB, for an act. & The upload is refused. The console refuses a larger file before sending it. \\
\bottomrule
\end{tabularx}
\end{table}

\paragraph{The team budget} is the hard limit of the whole installation, and it is not the
system's to raise.

\begin{lessonbox}{Observed, 27 September 2026}
The gateway refused every call with \enquote{Team budget exceeded: 10.18 >= 10.00 dollars}
(\file{docs/markdown/en/operations.md}). The budget of the key was ten dollars then. The product names
the condition by one code everywhere, \code{AI_BUDGET_EXHAUSTED}; a job step that meets it
fails at once instead of retrying, and the corpus commands take by default only what is still
missing, so a run interrupted by it is finished by running the command again.
\end{lessonbox}

\paragraph{The daily budget} keeps one account from draining the shared key. It is counted
per account and per UTC day from the recorded cost of that account's model calls
(\code{budget_status} in \file{assistant/usage.py}). An account's own amount, when an
administrator has set one, replaces the installation's default; administrators are unlimited
unless given an amount. What counts is every call recorded against the account: chat turns,
the steps of its reports and work products, prompt test runs, the model's part in a citation
check, the embedding of its uploads, and the readings an administrator starts from an act's
page. Shared work on the corpus is charged to nobody: the extraction of obligations inside a
report, the readings an operator starts from the command line, ingestion and re-embedding.

Three behaviours at the limit are deliberate. A report is refused only at its start and when
its scope is approved; a run that has started is never stopped halfway, so an account can end
a day above its budget by the cost of one run. A citation check on an account with nothing
left does not fail: it judges headings by word overlap instead of asking the model. And the
indexing of an upload does not fail either: the \code{documents.index} job queues itself
again for 00:05~UTC, five minutes after the reset (\code{RESUME_AFTER_RESET} in
\file{companies/indexing.py}), and its result on the jobs page says that it is waiting for
the daily budget.

\paragraph{Quotas and the budget together.} \file{docs/markdown/en/operations.md} works out the worst case
the defaults allow: one account that fills its quota of 1000 files and 2~GiB with text causes
embedding costs of roughly \$65 to \$195 at the price of the embedding model --- but never
more than its daily budget in one day, because indexing waits for the reset. An account that
uploads and deletes in a loop spends its daily budget every day and no more. The usage page of
the console shows it; the remedies are a lower budget for the account, a lower quota, or
deactivating the account.

\paragraph{Sign-in attempts} are limited too, although they cost nothing: ten per account in
ten minutes (\code{EMAIL_LIMIT}, \code{EMAIL_WINDOW_S} in \file{users/api/auth_views.py}) and
\code{LOGIN_RATE_PER_IP}, 60 by default, per client address in a minute. On the live host
most visitors arrive under one address and share that second window (\cref{ch:security}). The
chat allows
\code{RATE_LIMIT_TURNS}, 12, turns per user in a minute.


\section{The scheduled work}\label{sec:ops:schedule}

There is no cron container. The worker loop itself, between jobs and at most once every
60~s, looks at a small table (\code{SCHEDULE} in \file{jobs/schedule.py},
\code{SCHEDULE_EVERY_S} in \file{jobs/management/commands/run_worker.py}) and queues a kind when
the newest job of that kind --- whatever its status, and whoever started it --- was created
longer ago than the kind's interval. A kind that is still queued or running is never queued
twice. Two kinds are scheduled, both at priority $-5$, below everything else:

\begin{itemize}
  \item \code{corpus.refresh_all}, every seven days (\cref{sec:ops:corpus}). A check of all
    sources an administrator starts by hand restarts the seven days.
  \item \code{ops.maintenance}, every day (\cref{tab:ops:maintenance}).
\end{itemize}

\begin{table}[H]
\caption{What the daily maintenance job removes (\file{ops/maintenance.py}). The job's result
on the jobs page gives the count of each.}\label{tab:ops:maintenance}
\begin{tabularx}{\textwidth}{@{}L{38mm}YL{45mm}@{}}
\toprule
\hd{What} & \hd{Rule} & \hd{Constant} \\
\midrule
Expired sessions & Sessions past their expiry. & --- (a session lives \code{SESSION_DAYS}, 30) \\
Old jobs and their steps & Finished, failed or cancelled more than 90 days ago --- except a job a report points at, which is kept as long as the report. For a report made in two parts that is the job of the second part; the first part's is purged like any other (\cref{ch:arch}). & \code{JOB_RETENTION_DAYS} \\
Old records of model calls & Older than 400 days. The record of a chat call holds what was said, so this is also the retention of that copy of chat content. & \code{TELEMETRY_RETENTION_DAYS} \\
Orphaned files & Files in the documents folder that no document row names and that are older than six hours: what a crash between writing the file and the row leaves behind. & \code{ORPHAN_GRACE_HOURS} \\
Expired rate-limit windows & Counting windows of the rate limits whose time is over. & --- \\
Chat uploads nobody used & Documents with no company, older than seven days, that no chat message names: a file dropped into the chat and never sent. At most 500 a run. & \code{ORPHAN_UPLOAD_DAYS}, \code{ORPHAN_UPLOAD_BATCH} \\
\bottomrule
\end{tabularx}
\end{table}

The job has one attempt; a run that fails is tried again by the next day's schedule.


\section{Health, logs and troubleshooting}\label{sec:ops:health}

\paragraph{The health probe.} \route{GET}{/api/health} (\code{health} in
\file{ops/api/views.py}) is public on purpose --- it carries no data, and a probe that needs
credentials is a probe nobody wires up. It tests readiness: the database must answer a
trivial query, and the uploads folder must be writable.

\begin{codeblock}
200  {"status": "ok", "service": "legalops-django", "database": "ok", "uploads": "ok"}
503  {"status": "error", "service": "legalops-django", "reason": "database",
      "database": "error", "uploads": "ok"}
\end{codeblock}

The answer names which check failed and never why: an exception's text would hand a stranger
host names and paths, so the reason is in the log. The assistant service has a probe of its
own, \code{/health} on its port inside the isolated network, which answers whether a model is
configured, which providers have credentials, and how many chat turns are running; it is not
published through the endpoint. The worker loop touches a heartbeat file
(\file{/tmp/worker-heartbeat}) on every turn and during long jobs; nothing reads it today.
The plainest sign of a live worker is a job that leaves the queue.

\paragraph{Logs.} Every container writes to its standard output; \texttt{docker compose logs}
reads them. Neither the compose file nor the Docker daemon rotates these logs, so they grow
until the container is removed, which every update does. The endpoint writes one JSON line a
request with the status and the duration. The API logs at level \code{INFO} in the format set
in \file{settings.py}; in production gunicorn adds an access line. The worker container's
log is the loop's log, since the container only runs the loop. PostgreSQL logs statements
slower than two seconds (\code{log_min_duration_statement=2000}).

\begin{codeblock}
docker compose logs -f --tail 100 legalops-django     # follow the API
docker compose logs -f legalops-worker                # follow the job loop
\end{codeblock}

For anything that involves a model, the console is the better log: the jobs page shows each
step of a job with its error, and the usage page shows every recorded model call with its
cost (\cref{ch:console}).

\paragraph{Common failures.} \Cref{tab:ops:trouble} lists the failures an operator is likely
to meet, most of them from \file{docs/markdown/en/operations.md}, each checked against the code that
produces the message.

\begin{xltabular}{\textwidth}{@{}L{43mm}YL{42mm}@{}}
\caption{Common failures, what each means and what to do}\label{tab:ops:trouble}\\
\toprule
\hd{What is seen} & \hd{What it means} & \hd{What to do} \\
\midrule
\endfirsthead
\toprule
\hd{What is seen} & \hd{What it means} & \hd{What to do} \\
\midrule
\endhead
\bottomrule
\endlastfoot
502 from the endpoint & The service behind that path is down or still starting --- during an update, until it is back. & \texttt{docker compose ps}, then the logs of the service. \\
403 with an empty body on every request from another site & The endpoint refuses a request whose \code{Origin} header does not end with the host it was sent to. & Expected. The proxy in front has to pass on the \code{Host} header the browser sent. \\
The update script ends with a migration error & On a first installation, the database was still loading the bootstrap dump: harmless. Otherwise a migration failed and the stack runs on the old schema. & Run \texttt{manage.py migrate} again in the API's container. \\
A job step fails with \enquote{Prompt '…' has no active version} & The database did not start from the bootstrap dump and \code{seed_prompts} never ran. & Run \code{seed_prompts}. \\
Jobs stay queued & No worker loop is running: the worker container is restarting, or the API's container is down. & The worker's log; \file{/tmp/worker.pid} in the API's container. \\
Health answers 503 with reason \code{database} & PostgreSQL is down, the password in \code{DATABASE_URL} is not the role's, or every connection is taken (\enquote{too many clients already} in the API's log). & The database container's log; taken connections are freed by restarting the API's container. \\
Health answers 503 with reason \code{uploads}; uploads fail with \enquote{Permission denied} & Files in \file{_DATA/django} belong to root. & Give the folder to user 1000; the update script does. \\
PostgreSQL restarts, complaining of the data directory's ownership or permissions & \file{_DATA/postgres} is not owned by user 1000 with mode 700. & Correct the owner and the mode. \\
The chat answers 503 & No gateway key, or no active version of \prompt{chat.assistant}. & Set \code{AI_GATEWAY_KEY}; run \code{seed_prompts}. \\
\code{AI_BUDGET_EXHAUSTED} in the chat; job steps fail with status 402 & The gateway's team budget is spent. & Have the budget raised at the gateway; then run the interrupted work again. \\
\code{BUDGET_EXCEEDED} (429) & The user's own daily budget is spent. & Wait for 00:00~UTC, or raise the account's budget on the users page. \\
The chat says the assistant is restarting (\code{SHUTTING_DOWN}) & An update is restarting the assistant service. & Wait a few seconds. \\
An act fails with \enquote{Cannot reach <host> … is it on the egress allowlist?} & The exit does not list the host, or the API's container has no pin for it. & \Cref{sec:ops:secrets}. \\
Sign-in answers 429 \code{RATE_LIMITED} & Ten attempts for the account in ten minutes, or the per-address limit. & Wait as long as the \code{Retry-After} header says. \\
A chat answer arrives all at once, or is cut off & The proxy in front buffers the stream, or its read timeout is too short. & \Cref{sec:ops:runs}. \\
A document stays \ui{Teksto nėra} (\enquote{no text}) & A scanned PDF without a text layer; the system has no OCR. & Upload a PDF with text, or a DOCX. \\
A report on a library too large to be sent whole finds no evidence anywhere; corpus search finds only exact words & The corpus was never embedded. & The one re-embedding (\cref{sec:ops:install}). \\
A document's indexing job says it is waiting for the daily budget & The uploader's budget is spent; the job returns at 00:05~UTC. & Nothing, or raise the budget. \\
A changed \file{endpoint/Caddyfile} or \file{exit/Caddyfile} has no effect & The container still runs the old file. & Restart that container. \\
\end{xltabular}


\section{The sample company script}\label{sec:ops:demo}

\file{scripts/demo_1009.py} builds the course's sample company through the public API, as a
user would through the console, so it doubles as a check that a new installation works from
sign-in to a queued report. It signs in with the credentials given in \code{ADMIN_EMAIL} and
\code{ADMIN_PASSWORD}, then:

\begin{enumerate}
  \item finds the company \enquote{1009 LTD (demo)} among the account's own companies, or
    creates it with its profile --- a Lithuanian crypto-asset service provider, a private
    limited company (\lt{UAB}) of 40 staff --- and records its MiCA licence as a credential
    of status \code{claimed}, with the note that the policies name it without a number;
  \item uploads every Markdown file of the sample's folder that the company does not have
    yet --- the repository's own copy,
    \file{django/compliance/prompts/evals/data/1009-ltd/documents}, or the folder
    \code{MATERIALS} names --- as a document of kind \code{policy} (the sample has 26: 13 on
    information security and 13 on financial technology), and waits up to 15 minutes for the
    worker to index them;
  \item with \code{--analyse}, asks for the cost estimate of a report limited to seven acts
    --- those with the short titles MiCA, DORA, TFR, AMLR, PPTFPĮ, KRĮ and GDPR --- prints
    it, and starts that report in Lithuanian.
\end{enumerate}

The script can be run again safely: it reuses the company and skips files already uploaded.
It needs only Python with the \code{requests} library, which the API's image has; its header
gives the command that runs it in that image. The policies' default folder is
the repository's own copy of the sample: the 26 files under
\file{prompts/evals/data/1009-ltd/documents}. Indexing
costs cents; a report costs what the estimate prints, and both are charged to the account the
script signs in with (\cref{ch:report} gives measured figures for a report on this company).


\section{Secrets and the egress allowlist, in brief}\label{sec:ops:secrets}

\Cref{ch:security} explains why the secrets and the allowlist are arranged as they are. What
an operator does with them is short.

\paragraph{Rotating a secret.} Compose reads \file{.env} only when it creates a container, so
every rotation ends with \texttt{docker compose up -d}, which recreates the containers whose
values changed.

\begin{table}[H]
\caption{The secrets, where each is kept, and what a rotation does}\label{tab:ops:secrets}
\begin{tabularx}{\textwidth}{@{}L{42mm}L{30mm}Y@{}}
\toprule
\hd{Secret} & \hd{Kept in} & \hd{Rotation} \\
\midrule
\code{AI_GATEWAY_KEY} & \file{.env} & Edited by hand after the gateway issued a new key; the assistant service is recreated and uses it. \\
\code{ASSISTANT_INTERNAL_SECRET} & \file{.env} & A new random value; the API and the assistant service are recreated together, and internal calls fail for the seconds in between. \\
\code{DJANGO_SECRET_KEY} & \file{.env} & A new random value. Users notice nothing: sessions are bearer tokens stored as hashes, and nothing is encrypted with this key. \\
\code{DBGATE_PASSWORD} and \code{DBGATE_AUTH_HEADER} & \file{.env} & Both lines deleted; the update script mints a new pair on its next run. \\
The database role's password & \file{docker-compose.yml} & Changed in PostgreSQL with \texttt{ALTER ROLE} and in the three places the compose file writes it: \code{DATABASE_URL}, \code{POSTGRES_PASSWORD} and \code{PASSWORD_con1} of the database tool, which logs in as the same role. \file{docs/markdown/en/operations.md} gives the commands, which pass the new value through the environment and never on a command line. \\
An administrator's password & The database, hashed & Changed in the console; after a lock-out, \texttt{create\_admin -{}-reset} (\cref{sec:ops:install}). \\
\bottomrule
\end{tabularx}
\end{table}

\paragraph{Letting a new host out.} A container on the isolated network reaches an outside
host only when two things are in place: a block in \file{exit/Caddyfile} that matches the
host's name in the TLS handshake and forwards to port 443 of that host, and a line under
\code{extra_hosts} of the consuming service in \file{docker-compose.yml} that points the name
at the exit's address, 172.28.10.30. Then \texttt{docker compose up -d} recreates the consumer,
and the exit is restarted because only its file changed. Today the exit lists four names: the
gateway, and the three sources of the acts. An external MCP server registered in the console
needs its host added the same way, or its connection test fails. Two sources are absent on
purpose, by \file{docs/markdown/en/operations.md}: \code{eur-lex.europa.eu} and \code{e-seimas.lrs.lt}
answer unattended requests with a browser challenge that a fetcher cannot pass.

\paragraph{What is left to the operator.} \file{docs/markdown/en/operations.md} ends with a checklist of
what no code covers --- among them the client address the outer proxy passes on, the hosting's
own tools that can read the project folder, the switch to
production mode and the database tool. \Cref{ch:security} takes those up as residual risks.
